On construit un modèle spectral canonique du facteur résiduel
$\Rres(s) = \zeta(s) / \big(\zetaplus(s)\zetaminus(s)\big)$ de la fonction
zêta de Riemann, à partir de la dynamique de persistance PT. Trois acquis
principaux sont établis.
\begin{enumerate}[nosep,leftmargin=*]
\item Une expression close des amplitudes spectrales
$\kappap(s) = 1 - (1 - p^{-s})\exp(\Ap\,p^{-s})$ telle que
$\det(I - \DRop(s))^{-1} = \Rres(s)$, vérifiée à $10^{-12}$ dans
$\Reop(s) > 1$ ; $\DRop$ est l'opérateur diagonal sur $\ell^{2}(\Primes)$
associé.
\item Une formule de trace analytique exacte
$-d \log \Rres/ds = \sum_p (\log p)\,[1/(p^s - 1) + \Ap\,p^{-s}]$, vérifiée
à $10^{-26}$, qui se lit comme une formule de trace de Berry--Keating sur
les longueurs d'orbites $\ell_p = \log p$.
\item Une trace régularisée opératorielle
$\Treg(s) = \Tr_{\mathrm{reg}}\!\big[(s - i\,\HPTBK)^{-1}\,\DRop(s)\big]$
dont les pôles distributionnels sur la ligne critique sont identifiés aux
zéros non triviaux de $\Rres(s)$. Sur $t \in [10, 200]$, la formule brute
capture 74 zéros sur 79 à $|\Delta| < 1.5$ (sur $t \in [10, 70]$ : $17/17$),
avec saturation au-delà de $P_{\max} = 3\,000$. La régularisation explicite
$R_1$ par lissage gaussien ($\eps = 0.2$) améliore substantiellement la
précision~: \textbf{$75/79$ zéros à $|\Delta| < 0.5$ et $47/79$ à
$|\Delta| < 0.1$} sur le même domaine, avec une médiane $|\Delta| = 0.080$
contre $0.52$ pour la formule brute. Le test étendu à $t \in [10, 500]$
($269$ zéros) confirme la robustesse~: \textbf{$267/269$ à $|\Delta| < 1.5$
($99.3\,\%$), $242/269$ à $|\Delta| < 0.5$ ($90.0\,\%$), $139/269$ à
$|\Delta| < 0.1$ ($51.7\,\%$)}, avec une capture par bande de 100 unités
qui reste $\ge 86.6\,\%$ à $|\Delta| < 0.5$ sur les cinq sous-bandes --- la
performance ne s'effondre pas au-delà de $T = 200$.
\end{enumerate}

La localisation $\Reop(s) = 1/2$ admet quatre lectures équivalentes
internes à PT~: cellule de Planck symplectique $u_\star\, p_\star = 2\pi$,
Jacobien de Haar multiplicatif $du = dp/p$, transformation unitaire
$e^{su/2}$, et condition aux bords antipériodique forcée par l'involution
$T_3$ du crible $\bmod\,3$. Le résidu numérique $1/8$ qui sépare la densité
semi-classique PT de la densité de Riemann--von Mangoldt admet une
sur-détermination cohomologique~: $1/8 = \sper^{2}/2 = \lambdaH$
(auto-couplage Higgs), $1/8 = c_1/N_{\mathrm{corners}}$ (première classe
de Chern du fibré spinoriel $\qplus/\qminus$), $1/8 = \sper/4$ (phase de
Berry par coin de la double barrière). L'identité $\sper^{2}/2 = \sper/4$
a pour seules solutions $\sper = 0$ et $\sper = 1/2$ ; la valeur $\sper = 0$
étant exclue par $\Tone$ (axiome PT), la cohérence des deux lectures
sur-détermine $\sper = 1/2$ sans le dériver indépendamment.

Une prédiction structurellement induite --- $\Delta_N(L \bmod q) = 1/(8q)$
pour les fonctions $L$ de Dirichlet primitives --- a été testée sur les
zéros pré-calculés de \texttt{LMFDB} ($q \in \{3,5,7,11\}$, $\sim 140$
zéros par caractère, $T \le 200$). Le test rejette la prédiction à
$\chi^{2}(\mathrm{PT}) = 7.49$ contre $\chi^{2}(H_0) = 0.023$ sur quatre
degrés de liberté. La calibration sur les $50$ premiers zéros de $\zeta$
montre que le résidu $1/8$ est une identité entre formes asymptotiques
(entre la densité semi-classique BK et la densité Riemann--von Mangoldt),
non un décalage mesurable au comptage des zéros. La discipline de
falsifiabilité du programme en sort renforcée et son domaine de validité
précisé~: le mécanisme PT décrit ici est spécifique à $\zeta$, pas à la
famille Dirichlet.

L'article distingue strictement \emph{comprendre} la structure des zéros
et \emph{prouver} leur localisation~: le programme atteint le premier
objectif au niveau opératoriel ; le second reste un problème ouvert au
sens classique, équivalent au prolongement analytique strict de $\Rres(s)$
hors du demi-plan de convergence.
